\subsection{Removing transition kernels over integral support}
\label{subsec:integral-support-reduction}

\begin{proof}[Proof of Theorem~\ref{thm:integral-support-reduction}]
Write $j_q:F(q)\to R=F(1)$ for the transition, available because
$q\leq1$. Multiplication with $F(1)$ and naturality show that
$I_q=\im j_q$ is an $R$--$R$-ideal. Functoriality gives
$I_q\subseteq I_p$ when $q\leq p$, and multiplicative naturality gives
$I_qI_p\subseteq I_{qp}$. The products are the ordinary products in
$R$ and are independent of representatives. The inclusion diagram
satisfies the unit, associativity and distributivity laws. Proper
ideals are not required to have their own units.

For the ideal $K$ of Theorem~\ref{thm:upper-unit-quotient}, the unit
context gives $K_q\subseteq\ker j_q$. Conversely, if $j_q(x)=0$,
then every product $uxv$, with $u\in F(p)$ and $v\in F(r)$,
transports to $j_p(u)j_q(x)j_r(v)=0$ in $F(1)$. Since
$pqr\vee1=1$, all defining context tests vanish and $x\in K_q$.
Thus $K_q=\ker j_q$, and the quotient $F(q)/K_q$ identifies with
$I_q$, compatibly with all transitions and products.

The matrix-fiber nilpotence bound and the projective-class, monoid
and group isomorphisms now follow from the preceding theorem.
Its full lifting proof preserves exact supports and includes zero
amplitudes at nonzero supports. This is the integral specialization
of the unrestricted reduction, not a claim that arbitrary
coefficient-kernel products vanish before transport.
\end{proof}
